\section{Conclusion}
\label{sec:conclusion}

Bounded addition and inversion can be implemented by a sparse resistive
network while retaining every source solution and every coordinate needed to
recover it. The resulting exact feasibility problem is $\ER$-complete even
under the simultaneous structural restrictions of
Theorem~\ref{thm:structural}. These restrictions therefore do not suffice to
make exact feasibility easier than general existential-real decision.

The same correspondence explains the algebraic and numerical behavior of
feasible voltage sets. Rational equivalence realizes precisely the compact
basic closed sets over $\Q$, while semialgebraic homeomorphism realizes all
compact semialgebraic topologies. Unique voltages can have arbitrarily high
algebraic degree, and infeasible networks can have doubly exponentially small
residuals. Exact decision, tolerance-based numerical evidence, and certificates
for a gap promise consequently require different guarantees.

For AC models, the transfer depends on the meaning of an angle limit.
Encoding zero cycle winding gives consistent real angles, and reactive
balance then gives both an exact equal-angle theorem and a quantitative
stability estimate. Principal-angle constraints alone allow winding states.
The explicit model definitions and cycle conditions identify where each
complexity and stability conclusion applies.
